\documentclass[12pt,a4paper]{article}
\usepackage[utf8]{vietnam}
\usepackage{amsmath,amssymb}
\usepackage{tabularx}
\usepackage{tikz}
\usepackage{xcolor}
\usepackage[top=0.6cm, bottom=0.6cm, left=0.6cm, right=0.6cm]{geometry}
\usepackage[dethi]{ex_test}

\SAOPTN{widthSA=2.8,heightSA=0.62}
\setlength{\parskip}{0pt}

% Dong ke cham mo chuan viet tay
\newcommand{\dongke}[1]{%
  \begingroup%
  \setlength{\parskip}{0pt}%
  \baselineskip=18.5pt%
  \vspace{1.0mm}%
  \foreach \i in {1,...,#1}{%
    \noindent{\color{gray!60}\dotfill}\par%
  }%
  \vspace{0.6mm}%
  \endgroup%
}

% Dinh nghia o bai tap tu luan trong bang
\newcommand{\oBaiTap}[3][6]{%
  \vspace{1.0mm}%
  \noindent\textbf{#2} #3\par%
  \dongke{#1}%
}

\begin{document}
\pagestyle{empty}
% =============================================================================
% ========================= MA DE 101 - MAT TRUOC (TRANG 1/2) =============
% =============================================================================
\noindent
\begin{tabularx}{\linewidth}{|X|c|}
\hline
\textbf{PHIẾU KIỂM TRA ĐỊNH KỲ TOÁN 6 - TUẦN 09 - BUỔI 02} & \textbf{Mã đề: 101 - Trang 1/2} \\
\hline
\end{tabularx}

\vspace{1.0mm}

\noindent
\begin{tabularx}{\linewidth}{|p{3.2cm}|X|}
\hline
\centering\textbf{Điểm số} & \centering\textbf{Lời nhận xét của thầy cô} \tabularnewline
\rule{0pt}{1.8cm} & \\
\hline
\end{tabularx}

\vspace{1.0mm}

\noindent\textbf{Câu 1 (0,25 điểm).} Cho $A = \{x \in \mathbb{N} \mid 15 \le x \le 45\}$. Số phần tử của $A$ là: \shortans{$31$}\par\vspace{0.25mm}
\noindent\textbf{Câu 2 (0,25 điểm).} Chữ số $x$ để số $\overline{37x}$ chia hết cho cả 2 và 5 là: \shortans{$0$}\par\vspace{0.25mm}
\noindent\textbf{Câu 3 (0,25 điểm).} Chữ số $x$ để số $\overline{2x5}$ chia hết cho 9 là: \shortans{$2$}\par\vspace{0.25mm}
\noindent\textbf{Câu 4 (0,25 điểm).} Trong các số $11, 15, 21, 27$, số nào là số nguyên tố? \shortans{$11$}\par\vspace{0.25mm}
\noindent\textbf{Câu 5 (0,25 điểm).} Tìm $ƯCLN(24, 36)$: \shortans{$12$}\par\vspace{0.25mm}
\noindent\textbf{Câu 6 (0,25 điểm).} Tìm $BCNN(10, 15)$: \shortans{$30$}\par\vspace{0.25mm}
\noindent\textbf{Câu 7 (0,25 điểm).} Phân tích số $60$ ra thừa số nguyên tố được: \shortans{$2^2 \cdot 3 \cdot 5$}\par\vspace{0.25mm}
\noindent\textbf{Câu 8 (0,25 điểm).} Chu vi hình chữ nhật có chiều dài $14\text{ cm}$, rộng $8\text{ cm}$ là: \shortans{$44\text{ cm}$}\par\vspace{0.25mm}
\noindent\textbf{Câu 9 (0,25 điểm).} Diện tích hình bình hành đáy $15\text{ cm}$, chiều cao $8\text{ cm}$ là: \shortans{$120\text{ cm}^2$}\par\vspace{0.25mm}
\noindent\textbf{Câu 10 (0,25 điểm).} Có $24$ bạn nam và $18$ bạn nữ, chia nhiều nhất thành bao nhiêu tổ đều nhau? \shortans{$6\text{ tổ}$}\par\vspace{0.25mm}

\vspace{0.3mm}
\noindent\textbf{Câu 11 (3,0 điểm). Thực hiện phép tính (hợp lý nếu có thể):}
\vspace{0.4mm}

\noindent
\begin{tabularx}{\linewidth}{|X|X|}
\hline
\oBaiTap[7]{a)}{$45 + 172 + 55 + 28$} & \oBaiTap[7]{b)}{$28 \cdot 76 + 28 \cdot 34 - 28 \cdot 10$} \\
\hline
\oBaiTap[7]{c)}{$4^5 : 4^3 + 3 \cdot 2^3 - 2026^0$} & \oBaiTap[7]{d)}{$150 : (19 - 4) + 2^3 \cdot 5$} \\
\hline
\oBaiTap[7]{e)}{$180 : [43 - (12 - 7)^2]$} & \oBaiTap[7]{f)}{$200 : \{70 - [36 - (11 - 7)^2]\}$} \\
\hline
\end{tabularx}

\newpage
% =============================================================================
% ========================= MA DE 101 - MAT SAU (TRANG 2/2) ==============
% =============================================================================
\noindent
\begin{tabularx}{\linewidth}{|X|}
\hline
\centering\textbf{TOÁN 6 - TUẦN 09 - BUỔI 02 | Mã đề: 101 - Trang 2/2} \\
\hline
\end{tabularx}

\vspace{1.2mm}

\noindent\textbf{Câu 12 (2,0 điểm). Tìm số tự nhiên $x$, biết:}
\vspace{0.5mm}

\noindent
\begin{tabularx}{\linewidth}{|X|X|}
\hline
\oBaiTap[5]{a)}{$(x - 15) : 4 = 12$} & \oBaiTap[5]{b)}{$3 \cdot 2^x - 12 = 36$} \\
\hline
\oBaiTap[7]{c)}{$48 \ \vdots \ x, 72 \ \vdots \ x$ và $10 < x < 20$} & \oBaiTap[7]{d)}{$x \ \vdots \ 20, x \ \vdots \ 25$ và $150 < x < 250$} \\
\hline
\end{tabularx}

\vspace{1.2mm}

\noindent
\begin{tabularx}{\linewidth}{|X|}
\hline
\vspace{0.5mm}
\textbf{Câu 13 (2,0 điểm).} Một trường THCS tổ chức cho học sinh khối 6 đi tham quan dã ngoại. Khi xếp học sinh lên xe 30 chỗ, xe 40 chỗ hay xe 45 chỗ đều vừa thừa ra 5 học sinh. Biết số học sinh khối 6 trong khoảng từ 700 đến 800 học sinh.\par
\textbf{a) (1,0 điểm)} (1,0 điểm) Tìm $BCNN(30, 40, 45)$.\par
\textbf{b) (1,0 điểm)} (1,0 điểm) Tính số học sinh khối 6 của trường đó.\par
\vspace{0.4mm}
\noindent\textit{Bài giải:}\dongke{11}
\\ \hline
\end{tabularx}

\vspace{1.2mm}

\noindent
\begin{tabularx}{\linewidth}{|X|}
\hline
\vspace{0.5mm}
\textbf{Câu 14 (0,5 điểm).} Cho tổng $S = 2 + 2^2 + 2^3 + \dots + 2^{60}$. Chứng minh rằng $S$ chia hết cho $7$.\par
\vspace{0.4mm}
\noindent\textit{Bài giải:}\dongke{8}
\\ \hline
\end{tabularx}

\newpage
% =============================================================================
% ========================= MA DE 102 - MAT TRUOC (TRANG 1/2) =============
% =============================================================================
\noindent
\begin{tabularx}{\linewidth}{|X|c|}
\hline
\textbf{PHIẾU KIỂM TRA ĐỊNH KỲ TOÁN 6 - TUẦN 09 - BUỔI 02} & \textbf{Mã đề: 102 - Trang 1/2} \\
\hline
\end{tabularx}

\vspace{1.0mm}

\noindent
\begin{tabularx}{\linewidth}{|p{3.2cm}|X|}
\hline
\centering\textbf{Điểm số} & \centering\textbf{Lời nhận xét của thầy cô} \tabularnewline
\rule{0pt}{1.8cm} & \\
\hline
\end{tabularx}

\vspace{1.0mm}

\noindent\textbf{Câu 1 (0,25 điểm).} Cho $B = \{x \in \mathbb{N} \mid 12 \le x \le 52\}$. Số phần tử của $B$ là: \shortans{$41$}\par\vspace{0.25mm}
\noindent\textbf{Câu 2 (0,25 điểm).} Chữ số $y$ để số $\overline{48y}$ chia hết cho cả 2 và 5 là: \shortans{$0$}\par\vspace{0.25mm}
\noindent\textbf{Câu 3 (0,25 điểm).} Chữ số $x$ để số $\overline{4x3}$ chia hết cho 9 là: \shortans{$2$}\par\vspace{0.25mm}
\noindent\textbf{Câu 4 (0,25 điểm).} Trong các số $13, 25, 33, 35$, số nào là số nguyên tố? \shortans{$13$}\par\vspace{0.25mm}
\noindent\textbf{Câu 5 (0,25 điểm).} Tìm $ƯCLN(18, 30)$: \shortans{$6$}\par\vspace{0.25mm}
\noindent\textbf{Câu 6 (0,25 điểm).} Tìm $BCNN(12, 16)$: \shortans{$48$}\par\vspace{0.25mm}
\noindent\textbf{Câu 7 (0,25 điểm).} Phân tích số $72$ ra thừa số nguyên tố được: \shortans{$2^3 \cdot 3^2$}\par\vspace{0.25mm}
\noindent\textbf{Câu 8 (0,25 điểm).} Chu vi hình chữ nhật có chiều dài $16\text{ cm}$, rộng $9\text{ cm}$ là: \shortans{$50\text{ cm}$}\par\vspace{0.25mm}
\noindent\textbf{Câu 9 (0,25 điểm).} Diện tích hình bình hành đáy $16\text{ cm}$, chiều cao $7\text{ cm}$ là: \shortans{$112\text{ cm}^2$}\par\vspace{0.25mm}
\noindent\textbf{Câu 10 (0,25 điểm).} Có $30$ bút và $45$ quyển vở, chia nhiều nhất thành bao nhiêu phần thưởng đều nhau? \shortans{$15\text{ phần}$}\par\vspace{0.25mm}

\vspace{0.3mm}
\noindent\textbf{Câu 11 (3,0 điểm). Thực hiện phép tính (hợp lý nếu có thể):}
\vspace{0.4mm}

\noindent
\begin{tabularx}{\linewidth}{|X|X|}
\hline
\oBaiTap[7]{a)}{$36 + 185 + 64 + 15$} & \oBaiTap[7]{b)}{$35 \cdot 68 + 35 \cdot 42 - 35 \cdot 10$} \\
\hline
\oBaiTap[7]{c)}{$5^5 : 5^3 + 2 \cdot 3^2 - 2025^0$} & \oBaiTap[7]{d)}{$160 : (24 - 8) + 3^3 \cdot 2$} \\
\hline
\oBaiTap[7]{e)}{$210 : [46 - (13 - 8)^2]$} & \oBaiTap[7]{f)}{$250 : \{70 - [36 - (12 - 8)^2]\}$} \\
\hline
\end{tabularx}

\newpage
% =============================================================================
% ========================= MA DE 102 - MAT SAU (TRANG 2/2) ==============
% =============================================================================
\noindent
\begin{tabularx}{\linewidth}{|X|}
\hline
\centering\textbf{TOÁN 6 - TUẦN 09 - BUỔI 02 | Mã đề: 102 - Trang 2/2} \\
\hline
\end{tabularx}

\vspace{1.2mm}

\noindent\textbf{Câu 12 (2,0 điểm). Tìm số tự nhiên $x$, biết:}
\vspace{0.5mm}

\noindent
\begin{tabularx}{\linewidth}{|X|X|}
\hline
\oBaiTap[5]{a)}{$(x - 18) : 5 = 14$} & \oBaiTap[5]{b)}{$4 \cdot 3^x - 15 = 93$} \\
\hline
\oBaiTap[7]{c)}{$60 \ \vdots \ x, 90 \ \vdots \ x$ và $10 < x < 20$} & \oBaiTap[7]{d)}{$x \ \vdots \ 15, x \ \vdots \ 25$ và $100 < x < 200$} \\
\hline
\end{tabularx}

\vspace{1.2mm}

\noindent
\begin{tabularx}{\linewidth}{|X|}
\hline
\vspace{0.5mm}
\textbf{Câu 13 (2,0 điểm).} Một trường THCS tổ chức cho học sinh khối 6 đi tham quan dã ngoại. Khi xếp học sinh lên xe 24 chỗ, xe 36 chỗ hay xe 40 chỗ đều vừa thừa ra 3 học sinh. Biết số học sinh khối 6 trong khoảng từ 700 đến 800 học sinh.\par
\textbf{a) (1,0 điểm)} (1,0 điểm) Tìm $BCNN(24, 36, 40)$.\par
\textbf{b) (1,0 điểm)} (1,0 điểm) Tính số học sinh khối 6 của trường đó.\par
\vspace{0.4mm}
\noindent\textit{Bài giải:}\dongke{11}
\\ \hline
\end{tabularx}

\vspace{1.2mm}

\noindent
\begin{tabularx}{\linewidth}{|X|}
\hline
\vspace{0.5mm}
\textbf{Câu 14 (0,5 điểm).} Cho tổng $A = 3 + 3^2 + 3^3 + \dots + 3^{60}$. Chứng minh rằng $A$ chia hết cho $13$.\par
\vspace{0.4mm}
\noindent\textit{Bài giải:}\dongke{8}
\\ \hline
\end{tabularx}

\newpage
% =============================================================================
% ========================= MA DE 103 - MAT TRUOC (TRANG 1/2) =============
% =============================================================================
\noindent
\begin{tabularx}{\linewidth}{|X|c|}
\hline
\textbf{PHIẾU KIỂM TRA ĐỊNH KỲ TOÁN 6 - TUẦN 09 - BUỔI 02} & \textbf{Mã đề: 103 - Trang 1/2} \\
\hline
\end{tabularx}

\vspace{1.0mm}

\noindent
\begin{tabularx}{\linewidth}{|p{3.2cm}|X|}
\hline
\centering\textbf{Điểm số} & \centering\textbf{Lời nhận xét của thầy cô} \tabularnewline
\rule{0pt}{1.8cm} & \\
\hline
\end{tabularx}

\vspace{1.0mm}

\noindent\textbf{Câu 1 (0,25 điểm).} Cho $C = \{x \in \mathbb{N} \mid 20 \le x \le 60\}$. Số phần tử của $C$ là: \shortans{$41$}\par\vspace{0.25mm}
\noindent\textbf{Câu 2 (0,25 điểm).} Chữ số $a$ để số $\overline{59a}$ chia hết cho cả 2 và 5 là: \shortans{$0$}\par\vspace{0.25mm}
\noindent\textbf{Câu 3 (0,25 điểm).} Chữ số $x$ để số $\overline{5x1}$ chia hết cho 9 là: \shortans{$3$}\par\vspace{0.25mm}
\noindent\textbf{Câu 4 (0,25 điểm).} Trong các số $17, 27, 35, 49$, số nào là số nguyên tố? \shortans{$17$}\par\vspace{0.25mm}
\noindent\textbf{Câu 5 (0,25 điểm).} Tìm $ƯCLN(28, 42)$: \shortans{$14$}\par\vspace{0.25mm}
\noindent\textbf{Câu 6 (0,25 điểm).} Tìm $BCNN(15, 25)$: \shortans{$75$}\par\vspace{0.25mm}
\noindent\textbf{Câu 7 (0,25 điểm).} Phân tích số $90$ ra thừa số nguyên tố được: \shortans{$2 \cdot 3^2 \cdot 5$}\par\vspace{0.25mm}
\noindent\textbf{Câu 8 (0,25 điểm).} Chu vi hình chữ nhật có chiều dài $15\text{ cm}$, rộng $7\text{ cm}$ là: \shortans{$44\text{ cm}$}\par\vspace{0.25mm}
\noindent\textbf{Câu 9 (0,25 điểm).} Diện tích hình bình hành đáy $18\text{ cm}$, chiều cao $6\text{ cm}$ là: \shortans{$108\text{ cm}^2$}\par\vspace{0.25mm}
\noindent\textbf{Câu 10 (0,25 điểm).} Cắt hai thanh gỗ dài $36\text{ cm}$ và $48\text{ cm}$ thành các đoạn dài nhất bằng nhau. Độ dài mỗi đoạn là: \shortans{$12\text{ cm}$}\par\vspace{0.25mm}

\vspace{0.3mm}
\noindent\textbf{Câu 11 (3,0 điểm). Thực hiện phép tính (hợp lý nếu có thể):}
\vspace{0.4mm}

\noindent
\begin{tabularx}{\linewidth}{|X|X|}
\hline
\oBaiTap[7]{a)}{$52 + 164 + 48 + 36$} & \oBaiTap[7]{b)}{$42 \cdot 57 + 42 \cdot 53 - 42 \cdot 10$} \\
\hline
\oBaiTap[7]{c)}{$3^6 : 3^3 + 4 \cdot 2^3 - 2024^0$} & \oBaiTap[7]{d)}{$140 : (21 - 7) + 4^2 \cdot 3$} \\
\hline
\oBaiTap[7]{e)}{$240 : [49 - (14 - 9)^2]$} & \oBaiTap[7]{f)}{$300 : \{70 - [36 - (13 - 9)^2]\}$} \\
\hline
\end{tabularx}

\newpage
% =============================================================================
% ========================= MA DE 103 - MAT SAU (TRANG 2/2) ==============
% =============================================================================
\noindent
\begin{tabularx}{\linewidth}{|X|}
\hline
\centering\textbf{TOÁN 6 - TUẦN 09 - BUỔI 02 | Mã đề: 103 - Trang 2/2} \\
\hline
\end{tabularx}

\vspace{1.2mm}

\noindent\textbf{Câu 12 (2,0 điểm). Tìm số tự nhiên $x$, biết:}
\vspace{0.5mm}

\noindent
\begin{tabularx}{\linewidth}{|X|X|}
\hline
\oBaiTap[5]{a)}{$(x - 16) : 3 = 15$} & \oBaiTap[5]{b)}{$5 \cdot 2^x - 17 = 63$} \\
\hline
\oBaiTap[7]{c)}{$56 \ \vdots \ x, 84 \ \vdots \ x$ và $10 < x < 20$} & \oBaiTap[7]{d)}{$x \ \vdots \ 20, x \ \vdots \ 30$ và $150 < x < 200$} \\
\hline
\end{tabularx}

\vspace{1.2mm}

\noindent
\begin{tabularx}{\linewidth}{|X|}
\hline
\vspace{0.5mm}
\textbf{Câu 13 (2,0 điểm).} Một trường THCS tổ chức cho học sinh khối 6 đi tham quan dã ngoại. Khi xếp học sinh lên xe 20 chỗ, xe 25 chỗ hay xe 30 chỗ đều vừa thừa ra 4 học sinh. Biết số học sinh khối 6 trong khoảng từ 550 đến 650 học sinh.\par
\textbf{a) (1,0 điểm)} (1,0 điểm) Tìm $BCNN(20, 25, 30)$.\par
\textbf{b) (1,0 điểm)} (1,0 điểm) Tính số học sinh khối 6 của trường đó.\par
\vspace{0.4mm}
\noindent\textit{Bài giải:}\dongke{11}
\\ \hline
\end{tabularx}

\vspace{1.2mm}

\noindent
\begin{tabularx}{\linewidth}{|X|}
\hline
\vspace{0.5mm}
\textbf{Câu 14 (0,5 điểm).} Cho tổng $B = 4 + 4^2 + 4^3 + \dots + 4^{60}$. Chứng minh rằng $B$ chia hết cho $21$.\par
\vspace{0.4mm}
\noindent\textit{Bài giải:}\dongke{8}
\\ \hline
\end{tabularx}

\newpage
% =============================================================================
% ========================= MA DE 104 - MAT TRUOC (TRANG 1/2) =============
% =============================================================================
\noindent
\begin{tabularx}{\linewidth}{|X|c|}
\hline
\textbf{PHIẾU KIỂM TRA ĐỊNH KỲ TOÁN 6 - TUẦN 09 - BUỔI 02} & \textbf{Mã đề: 104 - Trang 1/2} \\
\hline
\end{tabularx}

\vspace{1.0mm}

\noindent
\begin{tabularx}{\linewidth}{|p{3.2cm}|X|}
\hline
\centering\textbf{Điểm số} & \centering\textbf{Lời nhận xét của thầy cô} \tabularnewline
\rule{0pt}{1.8cm} & \\
\hline
\end{tabularx}

\vspace{1.0mm}

\noindent\textbf{Câu 1 (0,25 điểm).} Cho $D = \{x \in \mathbb{N} \mid 18 \le x \le 68\}$. Số phần tử của $D$ là: \shortans{$51$}\par\vspace{0.25mm}
\noindent\textbf{Câu 2 (0,25 điểm).} Chữ số $b$ để số $\overline{62b}$ chia hết cho cả 2 và 5 là: \shortans{$0$}\par\vspace{0.25mm}
\noindent\textbf{Câu 3 (0,25 điểm).} Chữ số $x$ để số $\overline{6x1}$ chia hết cho 9 là: \shortans{$2$}\par\vspace{0.25mm}
\noindent\textbf{Câu 4 (0,25 điểm).} Trong các số $19, 21, 39, 45$, số nào là số nguyên tố? \shortans{$19$}\par\vspace{0.25mm}
\noindent\textbf{Câu 5 (0,25 điểm).} Tìm $ƯCLN(32, 48)$: \shortans{$16$}\par\vspace{0.25mm}
\noindent\textbf{Câu 6 (0,25 điểm).} Tìm $BCNN(14, 21)$: \shortans{$42$}\par\vspace{0.25mm}
\noindent\textbf{Câu 7 (0,25 điểm).} Phân tích số $84$ ra thừa số nguyên tố được: \shortans{$2^2 \cdot 3 \cdot 7$}\par\vspace{0.25mm}
\noindent\textbf{Câu 8 (0,25 điểm).} Chu vi hình chữ nhật có chiều dài $18\text{ cm}$, rộng $8\text{ cm}$ là: \shortans{$52\text{ cm}$}\par\vspace{0.25mm}
\noindent\textbf{Câu 9 (0,25 điểm).} Diện tích hình bình hành đáy $14\text{ cm}$, chiều cao $9\text{ cm}$ là: \shortans{$126\text{ cm}^2$}\par\vspace{0.25mm}
\noindent\textbf{Câu 10 (0,25 điểm).} Có $28$ bạn nam và $42$ bạn nữ, chia nhiều nhất thành bao nhiêu nhóm đều nhau? \shortans{$14\text{ nhóm}$}\par\vspace{0.25mm}

\vspace{0.3mm}
\noindent\textbf{Câu 11 (3,0 điểm). Thực hiện phép tính (hợp lý nếu có thể):}
\vspace{0.4mm}

\noindent
\begin{tabularx}{\linewidth}{|X|X|}
\hline
\oBaiTap[7]{a)}{$63 + 148 + 37 + 52$} & \oBaiTap[7]{b)}{$26 \cdot 84 + 26 \cdot 26 - 26 \cdot 10$} \\
\hline
\oBaiTap[7]{c)}{$2^7 : 2^3 + 5 \cdot 3^2 - 2023^0$} & \oBaiTap[7]{d)}{$180 : (23 - 5) + 5^2 \cdot 3$} \\
\hline
\oBaiTap[7]{e)}{$270 : [52 - (15 - 10)^2]$} & \oBaiTap[7]{f)}{$350 : \{70 - [36 - (14 - 10)^2]\}$} \\
\hline
\end{tabularx}

\newpage
% =============================================================================
% ========================= MA DE 104 - MAT SAU (TRANG 2/2) ==============
% =============================================================================
\noindent
\begin{tabularx}{\linewidth}{|X|}
\hline
\centering\textbf{TOÁN 6 - TUẦN 09 - BUỔI 02 | Mã đề: 104 - Trang 2/2} \\
\hline
\end{tabularx}

\vspace{1.2mm}

\noindent\textbf{Câu 12 (2,0 điểm). Tìm số tự nhiên $x$, biết:}
\vspace{0.5mm}

\noindent
\begin{tabularx}{\linewidth}{|X|X|}
\hline
\oBaiTap[5]{a)}{$(x - 20) : 6 = 11$} & \oBaiTap[5]{b)}{$2 \cdot 5^x - 25 = 225$} \\
\hline
\oBaiTap[7]{c)}{$64 \ \vdots \ x, 96 \ \vdots \ x$ và $10 < x < 20$} & \oBaiTap[7]{d)}{$x \ \vdots \ 18, x \ \vdots \ 24$ và $100 < x < 200$} \\
\hline
\end{tabularx}

\vspace{1.2mm}

\noindent
\begin{tabularx}{\linewidth}{|X|}
\hline
\vspace{0.5mm}
\textbf{Câu 13 (2,0 điểm).} Một trường THCS tổ chức cho học sinh khối 6 đi tham quan dã ngoại. Khi xếp học sinh lên xe 18 chỗ, xe 24 chỗ hay xe 30 chỗ đều vừa thừa ra 6 học sinh. Biết số học sinh khối 6 trong khoảng từ 700 đến 800 học sinh.\par
\textbf{a) (1,0 điểm)} (1,0 điểm) Tìm $BCNN(18, 24, 30)$.\par
\textbf{b) (1,0 điểm)} (1,0 điểm) Tính số học sinh khối 6 của trường đó.\par
\vspace{0.4mm}
\noindent\textit{Bài giải:}\dongke{11}
\\ \hline
\end{tabularx}

\vspace{1.2mm}

\noindent
\begin{tabularx}{\linewidth}{|X|}
\hline
\vspace{0.5mm}
\textbf{Câu 14 (0,5 điểm).} Cho tổng $C = 5 + 5^2 + 5^3 + \dots + 5^{60}$. Chứng minh rằng $C$ chia hết cho $31$.\par
\vspace{0.4mm}
\noindent\textit{Bài giải:}\dongke{8}
\\ \hline
\end{tabularx}
\end{document}
